% \documentclass[aspectratio=169,notes]{beamer}
\documentclass[aspectratio=169]{beamer}
\usetheme[faculty=phil]{fibeamer}
\usepackage{polyglossia}
\setmainlanguage{english} %% main locale instead of `english`, you
%% can typeset the presentation in either Czech or Slovak,
%% respectively.
\setotherlanguages{russian} %% The additional keys allow
%%
%%   \begin{otherlanguage}{czech}   ... \end{otherlanguage}
%%   \begin{otherlanguage}{slovak}  ... \end{otherlanguage}
%%
%% These macros specify information about the presentation
\title[Artificial Intelligence Software (AIS)]{Artificial Intelligence Software, HW 2} %% that will be typeset on the
\subtitle{VibeCoding  
\\ ROS2 \\ \  } %% title page.
\author{Oleg Bulichev}
%% These additional packages are used within the document:
\usepackage{ragged2e}  % `\justifying` text
\usepackage{booktabs}  % Tables
\usepackage{tabularx}
\usepackage{tikz}      % Diagrams
\usetikzlibrary{calc, shapes, backgrounds}
\usetikzlibrary{decorations.pathreplacing,calligraphy,calc,graphs}
\usepackage{amsmath, amssymb}
\usepackage{url}       % `\url`s
\usepackage{listings}  % Code listings
% \usepackage{subfigure}
\usepackage{floatrow}
\usepackage{subcaption}
\usepackage{mathtools}
\usepackage{todonotes}
\usepackage{fontspec}
\usepackage{multicol}
\usepackage{pdfpages}
\usepackage{wrapfig}
\usepackage{qrcode}
\usepackage{animate}
\usepackage{booktabs}
\usepackage{multirow}
% \usepackage{graphicx}
\usepackage{colortbl}

\graphicspath{{resources/}}
\frenchspacing

\setbeamertemplate{caption}[numbered]
\usetikzlibrary{graphs}

% \usepackage[backend=biber,style=ieee,autocite=footnote]{biblatex}
% \addbibresource{biblio.bib}
% \DefineBibliographyStrings{english}{%
%   bibliography = {References},}

\newcommand{\oleg}[2][] {\todo[color=red, #1] {OLEG:\\ #2}}
\newcommand{\fbckg}[1]{\usebackgroundtemplate{\includegraphics[width=\paperwidth]{#1}}}%frame background

\usepackage[framemethod=TikZ]{mdframed}
\newcommand{\dbox}[1]{
\begin{mdframed}[roundcorner=3pt, backgroundcolor=yellow, linewidth=0]
\vspace{1mm}
{#1}
\vspace{1mm}
\end{mdframed}
}

\begin{document}
\setlength{\abovedisplayskip}{0pt}
\setlength{\belowdisplayskip}{0pt}
\setlength{\abovedisplayshortskip}{0pt}
\setlength{\belowdisplayshortskip}{0pt}

\fbckg{fibeamer/figs/title_page.png}
\frame[c]{\setcounter{framenumber}{0}
    \usebeamerfont{title}%
    \usebeamercolor[fg]{title}%
    \begin{minipage}[b][6.5\baselineskip][b]{\textwidth}%
        \textcolor{black}{\raggedright\inserttitle}
    \end{minipage}
    % \vskip-1.5\baselineskip

    \usebeamerfont{subtitle}%
    \usebeamercolor[fg]{framesubtitle}%
    \begin{minipage}[b][3\baselineskip][b]{\textwidth}
        \raggedright%
        \insertsubtitle%
    \end{minipage}
    \vskip.25\baselineskip
}
%   \frame[c]{\maketitle}

\fbckg{fibeamer/figs/common.png}

\note{\scriptsize \begin{itemize}
        \item \
    \end{itemize}}

\begin{frame}[t]{Task 1}
\vspace{-0.3cm}
\textbf{Short description}:
Pass the full <<Beginner: CLI tools>> \href{https://docs.ros.org/en/humble/Tutorials/Beginner-CLI-Tools/Configuring-ROS2-Environment.html}{ROS2 tutorial}

\textbf{Artifacts}: Report in pdf format with your own screenshots.

\end{frame}

\begin{frame}[t]{Task 2}
\vspace{-0.3cm}
\textbf{Short description}:
Configure a Dockerized ROS~2 environment to stream camera feed, detect contours, and handle interactive zoom (Z/X). Use \texttt{rosdep} for dependency management.

\textbf{Artifacts}: \texttt{Dockerfile}, \texttt{docker-compose.yml}, \texttt{src/} (your ROS~2 package), and a PDF report with screenshots and usage instructions.


\textbf{Requirements}
\footnotesize
\vspace*{-0.4cm}
\begin{multicols}{2}
\begin{itemize}
    \item \textbf{Device Mapping}: Configure \texttt{docker-compose.yml} to access the host camera (e.g. \texttt{/dev/video0}).
    \item \textbf{Image Source}: Use a standard ROS~2 driver (like \texttt{v4l2\_camera} or \texttt{usb\_cam}) to publish images. Check with \texttt{rqt}.
    \item \textbf{Package Management}: Use \texttt{rosdep} to install dependencies. Do not install ROS packages via \texttt{apt} in \texttt{Dockerfile}.
    \item \textbf{Processing Node}: Create a node to detect contours on the video stream and publish the result on other ros topic.
    \item \textbf{Interaction}: Implement Zoom In (key 'Z') and Zoom Out (key 'X') functionality for the processed stream.
\end{itemize}
\end{multicols}

\end{frame}

\begin{frame}[t]{Task 3 (1)}
\vspace{-0.3cm}
\textbf{Context}:
You have collected raw field data: High-frequency Magnetometer logs (nanosecond timestamps, noisy) and Low-frequency GPS logs (NMEA format).



\textbf{Short description}:
Create a production-ready package to fuse these streams and visualize magnetic anomalies.

\textbf{Functional Requirements}:
\begin{itemize}
    \item \textbf{Data Pipeline}: Parse heterogeneous formats, synchronize data streams by time, and filter outliers.
    \item \textbf{Visualization}: Generate an interactive \textbf{Plotly} heatmap.
    \item \textbf{UX}: Implement dynamic threshold sliders and make a possibility to add OpenStreetMap background overlay.
\end{itemize}
\end{frame}

\begin{frame}[t]{Task 3 (2)}
\vspace{-0.3cm}
\textbf{Methodology}:
\begin{itemize}
    \item \textbf{TDD}: Strict Test-Driven Development. Write tests first. Coverage $>80\%$.
    \item \textbf{OOP}: Design a modular architecture (Abstract Parsers, \texttt{SyncStrategy} pattern).
    \item Use VibeCoding \textbf{MCP} for code verification and testing.
\end{itemize}

\textbf{Environment}:
\begin{itemize}
    \item Use \textbf{uv, Pixi} or other package manager.
    \item Commit \texttt{uv.lock} or \texttt{pixi.lock} for absolute reproducibility.
\end{itemize}
\end{frame}

\begin{frame}[t]{Task 3 (3)}
\vspace{-0.3cm}
\textbf{Artifacts}:
Docs in PDF format with the following sections:
\begin{itemize}
    \item \textbf{Architecture}: UML Class diagram of your solution.
    \item \textbf{Metrics Analysis}: Justify your choices for synchronization windows and filtering thresholds. How do they affect the result?
    \item \textbf{DevEx}: Reflection on the TDD process and modern tooling (\texttt{uv}/\texttt{pixi}, MCP).
\end{itemize}

\textbf{Repository}:
Clean structure with \texttt{src/}, \texttt{tests/}, and config files.
\end{frame}

\fbckg{fibeamer/figs/last_page.png}
\frame[plain]{}

\end{document}